\documentclass{notaaula}

        \titulonota{Estática e equilíbrio}
        \creditos{Turma 1B · Física · Dezembro/2026 · Prof. Josué Cavalcante}

        \begin{document}
        \cabecalho

        \begin{sobre}
        Equilíbrio de forças, torque, centro de gravidade e aplicações em estruturas e máquinas simples. Habilidades em foco: EM13CNT101, EM13CNT301, EM13CNT304 e EM13CNT307.
        \end{sobre}

        \section{Condições de equilíbrio}

\begin{copiar}{Translação e rotação}
Um corpo rígido permanece em equilíbrio estático quando não acelera nem translada e não ganha
rotação. As duas condições são
\[
  \sum \vec F=\vec0,
  \qquad
  \sum \tau=0.
\]
A primeira equação equilibra forças; a segunda equilibra as tendências de giro.
\atencao{Resultante de forças nula não garante, sozinha, equilíbrio: ainda pode existir torque resultante.}
\end{copiar}

\section{Momento de uma força}

\begin{copiar}{Torque}
O módulo do torque de uma força em relação a um eixo é
\[
  \tau=F\,d_{\perp}=F\,r\sen\theta.
\]
$d_{\perp}$ é a distância perpendicular entre o eixo e a linha de ação da força. Escolha um
sentido de rotação como positivo e mantenha a convenção em todo o cálculo.
\simbolos{$\tau$ ($\mathrm{N\,m}$); $F$ (N); $d_{\perp}$ e $r$ (m).}
\diaadia{Uma maçaneta longe das dobradiças aumenta o braço de alavanca e reduz a força necessária.}
\end{copiar}

\begin{exemplo}[Gangorra em equilíbrio]
Uma criança de peso $300\un{N}$ senta a $2\un{m}$ do apoio. Do outro lado, a que distância deve
sentar uma criança de peso $500\un{N}$?
\[
  300\cdot2=500\cdot d
  \Rightarrow d=\resultado{\dec{1,2}\un{m}}.
\]
\end{exemplo}

\section{Centro de gravidade e estabilidade}

\begin{copiar}{Base de apoio}
O centro de gravidade é o ponto em que podemos representar a ação do peso total. Um corpo
apoiado tende a permanecer estável enquanto a vertical do centro de gravidade cruza a base de apoio.
\begin{itens}
\item Base mais larga tende a aumentar a estabilidade.
\item Centro de gravidade mais baixo tende a aumentar a estabilidade.
\item Se a vertical do centro de gravidade sai da base, aparece torque que favorece o tombamento.
\end{itens}
\end{copiar}

\section{Estruturas e engenharia}

\begin{copiar}{Leitura de uma viga}
Em uma viga, os apoios exercem reações que devem equilibrar as cargas. Um procedimento seguro é:
\begin{itens}
\item desenhar o diagrama de forças;
\item escolher eixos e um ponto para calcular torques;
\item aplicar $\sum F_x=0$, $\sum F_y=0$ e $\sum\tau=0$;
\item verificar unidades, sinais e se as reações obtidas são fisicamente possíveis.
\end{itens}
Triângulos e treliças distribuem esforços com rigidez; cabos trabalham bem sob tração, e pilares
precisam resistir à compressão e à flambagem.
\end{copiar}

\begin{exemplo}[Viga simplesmente apoiada]
Uma viga horizontal de $4\un{m}$ tem peso desprezível e recebe carga de $600\un{N}$ no centro.
Por simetria e equilíbrio, as reações verticais nos dois apoios são
\[
  R_A=R_B=\resultado{300\un{N}}.
\]
\end{exemplo}

\begin{dica}
Calcular torque em torno de um apoio elimina da equação a força aplicada exatamente nesse apoio,
pois seu braço de alavanca é zero.
\end{dica}

\begin{exercicios}
\nivel{1}{Conceitos}
\begin{questoes}
\item Escreva as duas condições de equilíbrio de um corpo rígido.
\item Calcule o torque de $20\un{N}$ aplicado perpendicularmente a $\dec{0,50}\un{m}$ do eixo.
\item O torque aumenta ou diminui quando a força se aproxima do eixo?
\item Cite duas medidas que aumentam a estabilidade de um objeto.
\nivel{2}{Aplicação}
\item Uma força de $40\un{N}$ atua a $\dec{0,30}\un{m}$ do eixo e faz $30^\circ$ com a barra. Ache $\tau$.
\item Em uma gangorra, $200\un{N}$ atuam a $3\un{m}$ do apoio. Qual força a $2\un{m}$ equilibra o sistema?
\item Uma viga de $6\un{m}$ recebe carga central de $800\un{N}$. Ache as reações em apoios simétricos.
\item Por que é mais fácil abrir uma porta pela maçaneta do que perto da dobradiça?
\nivel{3}{Síntese}
\item Uma escada encostada pode deslizar mesmo sem se mover inicialmente. Identifique as forças relevantes.
\item Explique o tombamento usando a vertical do centro de gravidade e a base de apoio.
\item Uma placa de peso $240\un{N}$ está presa por dois suportes verticais simétricos. Que força cada suporte exerce?
\item Proponha uma mudança para tornar uma estante alta mais segura contra tombamento.
\end{questoes}
\gabarito{2) $10\un{N\,m}$; 3) diminui; 5) $6\un{N\,m}$;
6) $300\un{N}$; 7) $400\un{N}$ em cada apoio; 11) $120\un{N}$ em cada suporte.}
\end{exercicios}


\end{document}
